% =========================================================
% --- Ajustes Finos do Sumário (Normas UEMA) ---
% Reproduz fielmente o modelo (modelos-uema/sumario-modelo.pdf):
%  - Entradas (número + título) em negrito, tamanho 12 (Times)
%  - Número na margem esquerda e título afastado 1.3cm
%  - Espaçamento 1.5 (≈20,7pt) uniforme entre TODAS as linhas
%  - Números de página SEM negrito
%  - REFERÊNCIAS em caixa alta, recuada; APÊNDICE/ANEXO recuados com traço "–"
% =========================================================

% --- 0. Espaçamento vertical: exatamente 20,7pt entre as entradas (como no modelo) ---
% O modelo tem exatamente 20,7pt (= \baselineskip de 12pt com espaçamento 1.5)
% entre TODAS as linhas, inclusive antes de REFERÊNCIAS/APÊNDICES/ANEXOS.
\setlength{\cftbeforechapterskip}{0pt}
\setlength{\cftbeforesectionskip}{0pt}
\setlength{\cftbeforesubsectionskip}{0pt}
\setlength{\cftbeforesubsubsectionskip}{0pt}
\setlength{\cftbeforeparagraphskip}{0pt}
\setlength{\cftbeforesubparagraphskip}{0pt}
\setlength{\cftbeforepartskip}{0pt}
\setlength{\cftbeforebookskip}{0pt}
\renewcommand{\insertchapterspace}{}
\setlength{\cftparskip}{0pt}
% Garante \baselineskip = 20,7pt dentro do sumário
\AtBeginDocument{\addtocontents{toc}{\protect\setlength{\protect\baselineskip}{20.7pt}}}
% Espaço entre o título "SUMÁRIO" e a primeira entrada (~41,4pt no modelo)
\makeatletter
\renewcommand{\aftertoctitle}{\par\nobreak\vskip 20.9pt}%
\makeatother

% --- 1. Pontos de preenchimento (espaçamento fino, como no modelo) ---
\renewcommand{\cftdotsep}{2}
\renewcommand{\cftchapterdotsep}{2}
\renewcommand{\cftsectiondotsep}{2}
\renewcommand{\cftsubsectiondotsep}{2}
% Os PONTOS de preenchimento NÃO são negritos (o título continua negrito).
% O padrão do memoir aplica \bfseries no leader; removemos isso.
\renewcommand{\cftchapterleader}{\cftdotfill{\cftchapterdotsep}}
\renewcommand{\cftsectionleader}{\cftdotfill{\cftsectiondotsep}}
\renewcommand{\cftsubsectionleader}{\cftdotfill{\cftsubsectiondotsep}}

% --- 2. Indentação: número na margem (x=85pt) e título afastado 1.3cm (x=121pt) ---
% O abntex2 (modo ABNT 6027:2012) seta cada \cftXnnumwidth a partir de
% \cftlastnumwidth. Definimos por último, de forma direta e via \AtBeginDocument,
% para vencer qualquer reconfiguração posterior da classe.
\makeatletter
\setlength{\cftlastnumwidth}{1.3cm}
\AtBeginDocument{%
  \setlength{\cftlastnumwidth}{1.3cm}%
  \setlength{\cftchapternumwidth}{1.3cm}%
  \setlength{\cftsectionnumwidth}{1.3cm}%
  \setlength{\cftsubsectionnumwidth}{1.3cm}%
  \setlength{\cftsubsubsectionnumwidth}{1.3cm}%
}
\makeatother

\setlength{\cftchapterindent}{0pt}
\setlength{\cftsectionindent}{0pt}
\setlength{\cftsubsectionindent}{0pt}

% --- 3. Títulos (número + texto) em negrito, como no modelo ---
\renewcommand{\cftchapterfont}{\bfseries}
\renewcommand{\cftsectionfont}{\bfseries}
\renewcommand{\cftsubsectionfont}{\bfseries}
\renewcommand{\cftsubsubsectionfont}{\bfseries}

% --- 4. Números de página SEM negrito (no modelo são normais) ---
\renewcommand{\cftchapterpagefont}{\normalfont}
\renewcommand{\cftsectionpagefont}{\normalfont}
\renewcommand{\cftsubsectionpagefont}{\normalfont}
\renewcommand{\cftsubsubsectionpagefont}{\normalfont}

% =========================================================
% --- 5. Apêndices/Anexos: traço curto "–" entre número e título ---
% \ifuemasumarioapendice é ligado apenas dentro de \l@appendix, de modo que o
% traço apareça só em "APÊNDICE A – TÍTULO" / "ANEXO A – TÍTULO", com pouco
% espaçamento ao redor do hífen (como no modelo).
% =========================================================
\makeatletter
\newif\ifuemasumarioapendice
\renewcommand{\cftchapteraftersnum}{%
  \ifuemasumarioapendice\space--\space\fi
}%
% Para apêndices/anexos a caixa do número não deve ter \hfil (senão o título
% fica empurrado para longe do "–"). Usamos uma caixa natural.
\let\uemaoldchapternumberlinebox\chapternumberlinebox
\def\chapternumberlinebox#1#2{%
  \ifuemasumarioapendice
    {\hb@xt@#1{#2}}%
  \else
    \uemaoldchapternumberlinebox{#1}{#2}%
  \fi
}
\let\uemaoldlappendix\l@appendix
\renewcommand{\l@appendix}[2]{%
  \uemasumarioapendicetrue
  \uemaoldlappendix{#1}{#2}%
  \uemasumarioapendicefalse
}%
\makeatother

% =========================================================
% --- 6. Recuos e CAIXA ALTA de REFERÊNCIAS, APÊNDICES e ANEXOS ---
% No modelo, essas entradas ficam recuadas ~1.3cm, alinhadas ao título (x=121pt),
% e "REFERÊNCIAS" aparece em caixa alta.
% =========================================================

% 6a. Recuo de 1.3cm à esquerda de REFERÊNCIAS, APÊNDICES e ANEXOS.
%     Não usamos \cftsetindents global (é persistente e afetaria os capítulos).
%     Para REFERÊNCIAS ligamos um flag enquanto sua \contentsline é processada;
%     para APÊNDICES/ANEXOS o recuo é fixo dentro do próprio \l@appendix.
\makeatletter
\newif\ifuemaindentref
% Reindenta \l@chapter apenas quando o flag de referências está ativo
\let\uemaoldlchapter\l@chapter
\renewcommand{\l@chapter}[2]{%
  \begingroup
    \ifuemaindentref\setlength{\cftchapterindent}{1.3cm}\fi
    \uemaoldlchapter{#1}{#2}%
  \endgroup
  \global\uemaindentreffalse
}
% Reindenta \l@appendix (apêndices e anexos)
\let\uemaoldlappendixindent\l@appendix
\renewcommand{\l@appendix}[2]{%
  \begingroup
    \setlength{\cftchapterindent}{1.3cm}%
    \uemaoldlappendixindent{#1}{#2}%
  \endgroup
}
\makeatother

% 6b. REFERÊNCIAS: título em CAIXA ALTA e marcador de recuo no sumário.
\makeatletter
\renewcommand{\bibname}{REFERÊNCIAS}%
\renewcommand{\bibsection}{%
  \chapter*{\bibname}%
  \bibmark
  \ifnobibintoc\else
    \phantomsection
    \addtocontents{toc}{\protect\uemaindentreftrue}%
    \addcontentsline{toc}{chapter}{\texorpdfstring{\MakeTextUppercase{\bibname}}{\bibname}}%
    \addtocontents{toc}{\protect\uemaindentreffalse}%
  \fi
  \prebibhook
}
\makeatother
